% all_figures.tex  — DSFB Oil & Gas figure booklet
% Compile from the figures/ directory with:
%   pdflatex -interaction=nonstopmode all_figures.tex
\documentclass[12pt,a4paper]{article}
\usepackage[margin=2cm]{geometry}
\usepackage{graphicx}
\usepackage{caption}
\usepackage{hyperref}
\usepackage{booktabs}
\usepackage{array}
\usepackage{xcolor}
\definecolor{dsfbblue}{HTML}{2166ac}

\hypersetup{
  colorlinks=true, linkcolor=dsfbblue, urlcolor=dsfbblue,
  pdfauthor={DSFB Oil\&Gas},
  pdftitle={DSFB Oil and Gas — Figure Booklet},
}

\setlength{\parindent}{0pt}
\setlength{\parskip}{4pt}

\title{\textbf{DSFB Oil \& Gas — Figure Booklet}\\[4pt]
  \large 20 Empirical Figures from Real-World Datasets Only}
\author{Generated by \texttt{cargo run --release -- generate-figures}}
\date{\today}

\begin{document}
\maketitle
\tableofcontents
\newpage

%─────────────────────────────────────────────────────────────────────────────
\section{Data Provenance}
\label{sec:provenance}

All figures in this booklet are derived from real-world industrial
datasets.  Figures~16 and~17 additionally display synthetic reference bars
(shown in grey) for cross-dataset comparison; all other figures contain no
synthetic or simulated data.

\begin{center}
\begin{tabular}{>{\bfseries}lllr}
\toprule
Dataset & Source & Licence & Steps \\
\midrule
Petrobras 3W v2.0.0 & \url{github.com/petrobras/3W} & CC BY 4.0 & 9\,087 \\
Equinor Volve 15/9-F-15 & \url{data.equinor.com} & Volve Data Licence V1.0 & 5\,326 \\
RPDBCS ESPset & \url{github.com/Rogerio-Chaves/RPDBCS} & MIT & 6\,032 \\
\bottomrule
\end{tabular}
\end{center}

\textbf{Important academic-integrity note.} The DSFB grammar engine receives
only the residual signal; fault labels are withheld during processing and
appear in figures solely as post-hoc annotation metadata.  No claim of fault
detection or classification is made.

\clearpage
%─────────────────────────────────────────────────────────────────────────────
\section{Petrobras 3W Results}
\label{sec:3w}

\subsection*{Figure 1 — Full residual trace}
\begin{figure}[ht]
  \centering \includegraphics[width=\textwidth]{fig_01_3w_residual_annotated.pdf}
  \caption{Petrobras 3W: P-MON-CKP residual (kPa) for all 9\,087 steps
    coloured by DSFB grammar token.  Bottom strip shows 3W event-class labels
    (0 = normal; 1--8 = various fault scenarios).  CC BY 4.0.}
\end{figure}

\subsection*{Figure 2 — Per-well token distribution}
\begin{figure}[ht]
  \centering \includegraphics[width=0.9\textwidth]{fig_02_3w_token_per_well.pdf}
  \caption{Stacked horizontal bar chart of DSFB grammar-token fractions per
    well episode.  Wells sorted by ascending Nominal fraction.  CC BY 4.0.}
\end{figure}

\subsection*{Figure 3 — Phase portrait ($\tilde{r}$, $\tilde{\delta}$)}
\begin{figure}[ht]
  \centering \includegraphics[width=0.7\textwidth]{fig_03_3w_phase_portrait.pdf}
  \caption{Scatter plot in normalised residual--drift space.  Dashed box marks
    the $\pm$1 admissibility boundary.  Point colour = grammar token.  CC BY 4.0.}
\end{figure}

\subsection*{Figure 4 — Episode-length ECDF}
\begin{figure}[ht]
  \centering \includegraphics[width=0.7\textwidth]{fig_04_3w_episode_ecdf.pdf}
  \caption{Empirical CDF of well-episode lengths on a log-x axis.
    Vertical lines mark the 50th and 90th percentiles.  CC BY 4.0.}
\end{figure}

\subsection*{Figure 5 — Observed vs expected pressure}
\begin{figure}[ht]
  \centering \includegraphics[width=0.7\textwidth]{fig_05_3w_observed_expected.pdf}
  \caption{Scatter of observed versus expected P-MON-CKP coloured by 3W event
    class.  The dashed line is the $y=x$ identity.  CC BY 4.0.}
\end{figure}

\clearpage
%─────────────────────────────────────────────────────────────────────────────
\section{Equinor Volve Results}
\label{sec:volve}

\subsection*{Figure 6 — TQA residual trace}
\begin{figure}[ht]
  \centering \includegraphics[width=\textwidth]{fig_06_volve_tqa_annotated.pdf}
  \caption{Equinor Volve: surface torque (TQA) residual (kNm) vs measured depth,
    grammar token shading applied.  Volve Data Licence V1.0.}
\end{figure}

\subsection*{Figure 7 — Drift and slew channels}
\begin{figure}[ht]
  \centering \includegraphics[width=\textwidth]{fig_07_volve_drift_slew.pdf}
  \caption{Two-panel plot of drift ($\delta$, top) and slew ($\sigma$, bottom)
    vs depth with $\pm$1 envelope bounds dashed.  Volve Data Licence V1.0.}
\end{figure}

\subsection*{Figure 8 — Token distribution}
\begin{figure}[ht]
  \centering \includegraphics[width=0.8\textwidth]{fig_08_volve_token_dist.pdf}
  \caption{Horizontal bar chart of DSFB token fractions for all 5\,326 Volve
    steps.  Volve Data Licence V1.0.}
\end{figure}

\subsection*{Figure 9 — Episode length histogram}
\begin{figure}[ht]
  \centering \includegraphics[width=0.7\textwidth]{fig_09_volve_episode_hist.pdf}
  \caption{Histogram of Volve contiguous-anomaly episode lengths (log--log axes).
    Volve Data Licence V1.0.}
\end{figure}

\subsection*{Figure 10 — Phase portrait}
\begin{figure}[ht]
  \centering \includegraphics[width=0.65\textwidth]{fig_10_volve_phase_portrait.pdf}
  \caption{Normalised ($\tilde{r}$, $\tilde{\delta}$) scatter for Volve,
    coloured by grammar token.  Volve Data Licence V1.0.}
\end{figure}

\clearpage
%─────────────────────────────────────────────────────────────────────────────
\section{RPDBCS ESP Results}
\label{sec:esp}

\subsection*{Figure 11 — Per-unit EnvViolation vs fault rate}
\begin{figure}[ht]
  \centering \includegraphics[width=0.7\textwidth]{fig_11_esp_per_unit_envviol.pdf}
  \caption{For each of the 11 ESP units: true fault-step fraction (x-axis,
    ground-truth label) vs.\ DSFB EnvViolation-token fraction (y-axis).
    The dashed diagonal is the $y=x$ reference.
    \textbf{DSFB labels were withheld; no detection claim is made.}
    RPDBCS ESPset, MIT Licence.}
\end{figure}

\subsection*{Figure 12 — Broadband RMS by fault label}
\begin{figure}[ht]
  \centering \includegraphics[width=0.8\textwidth]{fig_12_esp_rms_by_class.pdf}
  \caption{Notched box plots of broadband RMS by fault category.
    RPDBCS ESPset, MIT Licence.}
\end{figure}

\subsection*{Figure 13 — Token distribution}
\begin{figure}[ht]
  \centering \includegraphics[width=0.8\textwidth]{fig_13_esp_token_dist.pdf}
  \caption{Grammar-token fraction for all 6\,032 ESP snapshots.
    RPDBCS ESPset, MIT Licence.}
\end{figure}

\subsection*{Figure 14 — Residual traces for representative units}
\begin{figure}[ht]
  \centering \includegraphics[width=\textwidth]{fig_14_esp_residual_units.pdf}
  \caption{DSFB residual traces for units 0, 1, and 4, token-shaded.
    RPDBCS ESPset, MIT Licence.}
\end{figure}

\subsection*{Figure 15 — Phase portrait coloured by fault label}
\begin{figure}[ht]
  \centering \includegraphics[width=0.7\textwidth]{fig_15_esp_phase_portrait.pdf}
  \caption{($\tilde{r}$, $\tilde{\delta}$) space coloured by ground-truth fault
    label (label \emph{not} provided to DSFB engine).
    RPDBCS ESPset, MIT Licence.}
\end{figure}

\clearpage
%─────────────────────────────────────────────────────────────────────────────
\section{Cross-Dataset Analysis}
\label{sec:cross}

\subsection*{Figure 16 — Noise compression ratio}
\begin{figure}[ht]
  \centering \includegraphics[width=0.85\textwidth]{fig_16_cross_ncr_bar.pdf}
  \caption{Cross-dataset NCR ($K/E$).  Real datasets shown in blue/purple;
    synthetic reference values in grey.  Asterisk: ESP NCR is low because the
    ESPset is a snapshot corpus, not a continuous stream.}
\end{figure}

\subsection*{Figure 17 — Event detection rate}
\begin{figure}[ht]
  \centering \includegraphics[width=0.85\textwidth]{fig_17_cross_edr_bar.pdf}
  \caption{EDR (Nominal fraction) across real and synthetic datasets.}
\end{figure}

\subsection*{Figure 18 — Token heatmap}
\begin{figure}[ht]
  \centering \includegraphics[width=\textwidth]{fig_18_cross_token_heatmap.pdf}
  \caption{Dataset $\times$ token percentage heatmap for the three real datasets.}
\end{figure}

\subsection*{Figure 19 — Episode-length violin}
\begin{figure}[ht]
  \centering \includegraphics[width=0.7\textwidth]{fig_19_cross_episode_violin.pdf}
  \caption{Violin plots (log-transformed) of contiguous anomaly episode lengths
    across the three real datasets.}
\end{figure}

\subsection*{Figure 20 — Envelope utilisation}
\begin{figure}[ht]
  \centering \includegraphics[width=0.8\textwidth]{fig_20_cross_envelope_utilisation.pdf}
  \caption{Fraction of steps with $|\tilde{r}|>1$, $|\tilde{\delta}|>1$, or
    $|\tilde{\sigma}|>1$ per dataset, revealing how often each DSFB dimension
    is spatially stressed in practice.}
\end{figure}

\end{document}
